\section{西周纲要：周人怎样把胜利变成秩序}

西周最值得追问的，不是“周武王打败商纣王”这一句话，而是胜利之后的难题：渭河流域的周人，怎样面对东方人口更密、旧商传统更强的地区？怎样让同姓宗族、功臣、旧商贵族和边疆部族在同一套秩序里各得其所？西周的答案不是后世的郡县制，而是一张由宗法、封国、婚姻、礼器、朝聘和军事义务织成的网络。

这张网络让周的统治可以延伸得很远，却也把风险埋在其中。王室强时，它能给各地分配名分、资源和安全；王室弱时，封国和宗族同样拥有独立行动的空间。西周末年的继承危机和西部边防失守，正是在这套网络无法自我修复时发生的。

\section{时间线：从牧野到东迁}
\begin{center}
\begin{tikzpicture}[x=.021cm,y=1cm]
  \draw[line width=.8pt,color=timeline] (0,0) -- (275,0);
  \foreach \x in {0,6,69,205,275} {
    \draw[color=dynasty,line width=.8pt] (\x,-.11) -- (\x,.11);
  }
  \node[above,align=center,font=\small\color{ink}] at (0,0) {约前1046\\牧野};
  \node[below,align=center,font=\small\color{ink}] at (6,0) {约前1040\\东征};
  \node[above,align=center,font=\small\color{ink}] at (69,0) {前977\\昭王南征};
  \node[below,align=center,font=\small\color{ink}] at (205,0) {前841\\共和};
  \node[above,align=center,font=\small\color{ink}] at (275,0) {前771\\镐京陷落};
\end{tikzpicture}
\end{center}

\begin{sourcenote}[这条时间线有多可靠]
西周早期的绝对年代仍有断代争议：牧野、周公摄政和成周建设的年份不能像秦汉以后那样写得毫无保留。公元前841年共和元年则是传统连续纪年的稳定锚点。以下条目会把“约”“传统说法”“可靠纪年”分开，不用整齐的年份掩盖材料本身的粗细不一。
\end{sourcenote}
